% =========================================================
\chapter{Lecture 8 : Adversarial thinking I}
% =========================================================

\begin{rembox}[title=In a nutshell]
\textbf{What is this lecture about?} Until now we built defenses. Here we switch sides and think
like the attacker. We see \textbf{why} studying attacks matters, \textbf{how} attacks are developed
(the \textbf{attack engineering process}: the security engineering process of Lecture 1 read
backwards, each step being a place where the defender may have made a mistake), with real-world
examples (HSM key extraction, hacked cars, fake GSM antennas, Tesla key fobs, WEP, the
\texttt{sudo} bug), and finally how a defender can \textbf{reason systematically} about ``what can
go wrong'' with threat modelling methodologies (\textbf{attack trees}, \textbf{STRIDE},
\textbf{P.A.S.T.A.}, security cards). \\
\textbf{Why is it useful?} You cannot defend against attacks you cannot imagine. The three-step
classification (policy and threat model / design / implementation and operation) is a checklist to
analyse any system, and STRIDE is a classic exam question: given a threat, say which property it
breaks.
\end{rembox}
